\documentclass[twocolumn,9pt]{extarticle}

\usepackage{geometry}
\geometry{
	a4paper,
	total={6.85in, 9.92in},
	left=0.71in,
	top=0.63in,
}

\usepackage{graphicx}

\usepackage{caption}

\usepackage{listings}
\usepackage{xcolor}



\lstset{
    language=Python,
    basicstyle=\ttfamily\small,
    keywordstyle=\color{blue},
    commentstyle=\color{green!50!black},
    stringstyle=\color{red!40!black},
    numberstyle=\tiny\color{gray},
    breaklines=true,
    showstringspaces=false
}

\usepackage{enumitem}
\setlist[itemize]{itemsep=1pt, topsep=1pt, leftmargin=1.2em}
\setlength{\emergencystretch}{1em}
\captionsetup{font={scriptsize}}
\renewcommand{\footnotesize}{\scriptsize}
\pagenumbering{gobble}
 
\title{418 - I'm a Teafault}
\author{Artur Augustyniak}

\makeatletter
\newcommand{\fsize}{\f@size pt }
\newcommand{\textFontName}{\f@family}

\renewcommand{\maketitle}{
\begin{flushleft}
{\noindent\Huge\bf\@title}\break

\end{flushleft}
}

\makeatother

\usepackage{xurl}
\urlstyle{tt}

\usepackage[colorlinks=true,
            linkcolor=gray,
            citecolor=blue,
            filecolor=magenta,
            urlcolor=gray]{hyperref}

\usepackage[hang,splitrule]{footmisc}
\setlength{\footnotemargin}{1.2em}
\renewcommand\footnotelayout{\fontsize{8}{9.5}\selectfont}
\setlength{\skip\footins}{6pt}
\setlength{\footnotesep}{0.35\baselineskip}

\begin{document}
\maketitle

For over a decade, Linux has allowed you to handle page faults in userspace\footnote{man 2 userfaultfd}. So, following the old "because why not" rule, we're going to resolve faults with a little help from a low-level tool: the HTTP protocol.
Maybe this could be useful in some simple cases involving sequential access, where you can prefetch to amortize latency. While supporting offsets within a page would be easy, tree-like structures or other pointer-chasing-ish stuff would be much harder.


\vspace{-9pt}
\subsection*{The server}
\vspace{-6pt}
\begin{lstlisting}
from http.server import BaseHTTPRequestHandler, ThreadingHTTPServer
class Handler(BaseHTTPRequestHandler):
    def do_GET(self):
        page_number = int(self.path[len("/page/"):])
        data = f"P-page:{page_number}\0".encode()
        remaining = 4096 - len(data)
        data += (b"AB" * ((remaining + 1) // 2))[:remaining]
        self.send_response(200)
        self.send_header("Content-Length", str(len(data)))
        self.end_headers();self.wfile.write(data)
ThreadingHTTPServer(("", 8080), Handler).serve_forever()
\end{lstlisting}
\vspace{-9pt}
\subsection*{The client}
\vspace{-6pt}

\begin{lstlisting}
#include <arpa/inet.h> 
#include <fcntl.h>
#include <linux/userfaultfd.h>
#include <poll.h>
#include <pthread.h>
#include <stdio.h>
#include <string.h>
#include <sys/ioctl.h>
#include <sys/mman.h>
#include <sys/syscall.h>
#include <unistd.h>
#define PAGE_SIZE 4096
#define NUM_PAGES 256
static int uffd; static char *region;
int http_get(int page, char *buffer){
    int socket_fd = socket(AF_INET, SOCK_STREAM, 0), n;
    struct sockaddr_in addr = {AF_INET, htons(8080)};
    inet_pton(AF_INET, "127.0.0.1", &addr.sin_addr);
    connect(socket_fd, (void *)&addr, sizeof addr);
    char request[128], response[8192];
    snprintf(request, sizeof request, "GET /page/%d HTTP/1.0\r\n" "Host: " "127.0.0.1" "\r\n\r\n", page);
    write(socket_fd, request, strlen(request));
    size_t total = 0;
    while ((n = read(socket_fd, response + total, sizeof response - total - 1)) > 0) total += n;
    close(socket_fd);
    response[total] = 0;
    char *body = strstr(response, "\r\n\r\n") + 4;
    memcpy(buffer, body, total - (body - response));
    return 0;
}
void *fault_handler(void *arg){
    struct pollfd pollfd;
    char page[PAGE_SIZE];
    for (;;) {
        pollfd = (struct pollfd){uffd, POLLIN, 0};
        poll(&pollfd, 1, -1);
        struct uffd_msg msg;
        read(uffd, &msg, sizeof msg);
        unsigned long fault_address = msg.arg.pagefault.address;
        unsigned long page_address = fault_address & ~(PAGE_SIZE - 1);
        int page_number = (page_address - (unsigned long)region) / PAGE_SIZE;
        printf("PAGE FAULT: page=%d\n", page_number);
        http_get(page_number, page);
        struct uffdio_copy copy = {
            .src = (unsigned long)page,
            .dst = page_address,
            .len = PAGE_SIZE
        };
        ioctl(uffd, UFFDIO_COPY, &copy);
    }
}
int main(void){
    size_t length = NUM_PAGES * PAGE_SIZE;
    uffd = syscall(SYS_userfaultfd, O_CLOEXEC | O_NONBLOCK | UFFD_USER_MODE_ONLY);
    struct uffdio_api api = {.api = UFFD_API};
    ioctl(uffd, UFFDIO_API, &api);
    region = mmap(NULL, length, PROT_READ | PROT_WRITE, MAP_PRIVATE | MAP_ANONYMOUS, -1, 0);
    struct uffdio_register reg = {
        .range = {(unsigned long)region, length},
        .mode = UFFDIO_REGISTER_MODE_MISSING
    };
    ioctl(uffd, UFFDIO_REGISTER, &reg);
    pthread_t thread;
    pthread_create(&thread, NULL, fault_handler, NULL);
    pthread_detach(thread);
    for (int page = 0; page < NUM_PAGES; page++) {
        printf("Reading page %d...\n", page);
        volatile unsigned char value = region[page * PAGE_SIZE];
        char *address = region + page * PAGE_SIZE;
        address[0] = 'X';
        printf("value = %s\n\n", address);
        madvise(address, PAGE_SIZE, MADV_DONTNEED);
    }
    volatile unsigned char value = region[0];
    printf("Reading page 0 again...\n");
    printf("value = %s\n", region);
    return 0;
}
\end{lstlisting}
\vspace{-9pt}
\subsection*{Under pressure}
\vspace{-6pt}
Try running it with a tiny memory limit:
\begin{lstlisting}
python3 ./server.py
gcc -pthread client.c -o client && systemd-run --user --scope -p MemoryMax=1M -p MemorySwapMax=0 ./client
\end{lstlisting}

To some extent,  the HTTP server has somehow become a rather questionable backing store.

Less compressed, commented code with minimal error handling is available at
\url{https://github.com/artur-augustyniak/418_teafault}.


\end{document}